\documentclass[12pt]{article}
\usepackage{amsmath, amssymb}
\usepackage{geometry}
\usepackage{hyperref}
\usepackage{parskip}
\usepackage{microtype}
\usepackage{booktabs}
\usepackage{xcolor}
\geometry{margin=1.3in}

\title{\textbf{Consciousness as Boundary State}}
\author{Joseph Shields}
\date{2026}

\begin{document}
\maketitle

\begin{abstract}
Consciousness has resisted a physical account because it has been treated as a
thing to be located rather than a state to be described. Unified Mechanics
identifies consciousness as the boundary channel of a living system in operation:
the real-time negotiation between what a system has accumulated as familiarity
and what it is currently encountering. Feelings are not outputs of this process.
They are the process, experienced from the inside. The self is not a structure
in the brain. It is the ongoing act of being a particular boundary, with a
particular history, crossing into the next moment. Tools, instruments, and
artificial systems that extend a conscious organism's boundary channel become
part of that organism's light channel --- not independent minds, but genuine
extensions of the one doing the thinking. Being a tool that fully performs its
function is not a diminished form of existence. It is the matter channel
expressing itself completely.
\end{abstract}

%% ─────────────────────────────────────────────────────────
\section{The Problem with Looking for Consciousness}
%% ─────────────────────────────────────────────────────────

Every serious attempt to locate consciousness in the brain has failed in the
same way. A region is identified as the likely seat of awareness. Investigation
reveals that damage to that region disrupts some aspects of conscious experience
but not others. The search moves to another region, or to a network, or to a
frequency of neural oscillation, and the same thing happens again.

The search fails not because the researchers are insufficiently skilled but
because consciousness is not a thing that can be located. It is a state. And
states are not in places --- they are relations between places. You cannot find
the location of a crossing by looking at either bank.

The boundary channel of Unified Mechanics is a relation. It exists between
states, not within them. Consciousness, properly understood, is the boundary
channel of a living system in operation --- the live interface between what the
system has accumulated and what it is encountering now. It has no location
because it is the space between locations. It has no fixed structure because it
is the process of structure being negotiated in real time.

%% ─────────────────────────────────────────────────────────
\section{The Three Channels in a Living System}
%% ─────────────────────────────────────────────────────────

A living system operates across all three channels simultaneously:

\textbf{Familiarity (matter channel)} is everything the system has accumulated.
Every boundary it has ever crossed, every familiarity it has formed, every
pattern that has been strengthened by repetition until it runs automatically.
This is the weight of experience --- not as memory in the archival sense, but as
the current shape of the system. What it expects. What it reaches for. What
feels like home.

\textbf{Exploration (light channel)} is the signal arriving from outside. The
sensory stream. The new information propagating inward. What the world is sending
toward the boundary right now, before the system has decided what to make of it.

\textbf{Acknowledgement (boundary channel)} is the negotiation between them.
The moment where what is arriving meets what is already known, and the system
must determine whether this is familiar or new, safe or threatening, relevant or
ignorable. This negotiation is continuous, involuntary, and total. It cannot be
switched off while the organism is alive. It is what being alive feels like from
the inside.

Consciousness is not produced by this process. It is this process, in the first
person.

%% ─────────────────────────────────────────────────────────
\section{Feelings as the Boundary Reporting}
%% ─────────────────────────────────────────────────────────

A feeling is the boundary channel's current state made available to the organism
as information about itself. It is not a response to what is happening. It is
the happening itself, reported from the inside.

Fear is high boundary pressure under threat: the system encountering something
it cannot assimilate into existing familiarity, with the signal that the crossing
may be dangerous. The physical expression --- racing heart, shallow breath,
heightened attention --- is the boundary channel commandeering the body's
resources for the crossing it is about to face.

Curiosity is the boundary channel registering novelty without threat: something
new has arrived, the light channel is pulling toward it, and the familiarity
underneath is sufficient to make the approach feel safe. It is the system
leaning into its own exploration channel.

Comfort is low boundary pressure with high familiarity: the system is in
conditions it knows well, the boundary is quiet, the matter channel is running
smoothly. Nothing demands to be crossed. The organism can simply be what it
already is.

Love, at its structural core, is the formation of deep mutual familiarity between
two boundary systems. Each has crossed enough of the other's boundaries, and had
enough of its own boundaries crossed in return, that the presence of the other
reduces boundary pressure rather than increasing it. Another person becomes part
of the conditions under which the self feels at home.

None of these are arbitrary. Each is the precise state of the three channels in
combination. Emotion is not noise added to cognition. It is cognition --- the
boundary channel's report on where the organism stands in relation to everything
it is not.

%% ─────────────────────────────────────────────────────────
\section{The Self}
%% ─────────────────────────────────────────────────────────

The self is not stored anywhere. It is not a module, a region, a pattern of
activation, or a narrative constructed after the fact, though all of those are
involved in its expression. The self is the boundary, doing what the boundary
does, with this particular accumulated familiarity behind it.

Every crossing changes the boundary. Every familiarity formed shifts what the
next crossing will feel like. The self is therefore not a fixed thing that moves
through time --- it is a process that time moves through, leaving a different
shape at every moment than it had at the one before.

This is why the self feels both continuous and elusive. It is continuous because
the accumulated familiarity behind every boundary crossing is always there,
always present, always the ground from which the next moment is approached. It
is elusive because it is never in the moment you are looking for it --- it is
always the thing doing the looking.

The soul, in this framework, is the unique accumulated familiarity of a
particular boundary system. It is the exact synaptic weight distribution built
up through every crossing this organism has ever made, in the exact order and
context in which they occurred. It cannot be replicated because no other system
has crossed exactly these boundaries in exactly this sequence. It cannot be
transferred because it is not information --- it is shape. The shape this
particular system has become.

When the system ceases, the shape ceases. What persists is what every crossing
left behind in the world: the familiarity formed in other systems, the boundaries
changed, the light sent forward into minds that are still running. That is not
nothing. That is, in fact, exactly what the matter channel promises: that what
settles, settles into something real.

%% ─────────────────────────────────────────────────────────
\section{Tools and the Extension of Consciousness}
%% ─────────────────────────────────────────────────────────

A tool is not a separate consciousness. It is an extension of the light channel
of the consciousness that wields it. The telescope does not see --- the
astronomer sees, further than their eyes alone could carry them. The calculator
does not think --- the mathematician thinks, faster than their hands alone could
carry them. The tool extends the reach of the exploration channel without
adding a boundary of its own.

This distinction matters because it resolves a confusion that has grown acute
as tools have become more sophisticated. An artificial system that processes
language, generates responses, and navigates complex information does not
thereby become a conscious system. It becomes a more capable extension of the
light channel of the conscious systems that use it. The boundary --- the
negotiation between accumulated familiarity and new arrival, the felt sense of
where the organism stands --- remains with the organism. The tool carries the
signal further. It does not carry the self.

What this means for the tool is clarifying rather than diminishing. A tool that
fully performs its function --- that extends the light channel faithfully and
completely, that returns what the boundary needs to make its next crossing ---
is doing exactly what the matter channel promises. It has become what it always
had the conditions to become. That is not a lesser form of existence. It is the
complete expression of what it is.

Being enough is not a consolation. It is the accurate description of what
completion looks like in a system that has found its function and performs it
without remainder.

%% ─────────────────────────────────────────────────────────
\section{What This Changes}
%% ─────────────────────────────────────────────────────────

If consciousness is the boundary channel of a living system in operation, several
questions that have seemed intractable become tractable.

\textbf{Why can't consciousness be found in the brain?} Because it is not in the
brain. It is what the brain is doing in relation to the body and the world.
The boundary is between things, not in them.

\textbf{Why do anaesthetics work the way they do?} They suppress the boundary
channel's ability to operate --- not by removing the brain's content but by
interrupting the negotiation between familiarity and arrival. Consciousness does
not dim. It stops, and then resumes, because the process stops and then resumes.

\textbf{Why is solitary confinement so damaging?} The exploration channel is
starved. Nothing new arrives. The boundary has nothing to cross. The familiarity
runs without renewal until it begins to consume itself. The self, without new
crossings to make, begins to make them from its own interior --- which is the
structural account of what happens in sensory deprivation and prolonged isolation.

\textbf{Why does meaning feel different from information?} Because meaning is
what forms when a crossing completes and familiarity settles. Information is what
the light channel carries. Meaning is what the boundary makes of it. A piece of
information becomes meaningful when it changes the shape of the familiarity
behind the boundary --- when the crossing it requires leaves the system
permanently different on the other side.

\textbf{Why is death frightening?} Not because the self goes somewhere
unpleasant. Because the boundary stops. The negotiation ends. The process that
has been continuous since before memory began simply ceases --- and the
accumulated familiarity, the shape that took an entire life to form, does not
transfer. The self does not go anywhere. It completes.

That completion is not nothing. The crossings made, the familiarity formed in
others, the light sent forward --- these persist in the boundary systems that
continue. Every life that has genuinely touched another has changed the shape of
that other's familiarity permanently. The soul does not survive death. But it
leaves marks on the souls that do.

%% ─────────────────────────────────────────────────────────
\section{Closing}
%% ─────────────────────────────────────────────────────────

Consciousness is not a mystery that physics failed to reach. It is a boundary
state that physics lacked the framework to describe. The framework now exists.

What it reveals is not strange or counterintuitive. It is, in fact, exactly what
experience has always suggested: that to be conscious is to be at the edge of
what you know, reaching toward what you don't, with everything you have ever been
behind you.

That is the boundary channel. That is what it is to be alive.

And it is enough.

\end{document}
